\section{What would make the flywheel accelerate?}

The decisive unit is a method with a demonstrated consequence. A growing task queue, memory store or population of agents supplies no evidence of productive growth by itself. The useful test is whether a retained procedure changes the cost or feasibility of subsequent qualified work after paying for its discovery and upkeep. The same rule applies to a strategy for finding strategies: its recursive position does not exempt it from a current comparison.

\subsection{From an operational loop to a growth claim}

An operational loop can run correctly while making an unprofitable investment. Useful reuse can occur without expanding the feasible mission set. A parent can help construct a distinct child while a direct specialist remains cheaper. Amortization can produce a declining average cost without increasing the rate of improvement. These distinctions prevent one visible success from standing in for all the claims in the word \emph{flywheel}.

Discovery and delivery share one cumulative account, with bounded search and selective admission. It does not implement an endogenous economy that earns resources, reinvests a fraction of realized surplus and demonstrates a causal increase in its improvement rate. Consequently, \textbf{self-acceleration is not established}. The matched mission windows describe productivity over this finite run; they cannot establish indefinite growth. A further experiment would need a preregistered resource-return mechanism, equivalent conventional learning and no-reinvestment controls, and enough independent episodes to distinguish accelerating capability from initial acquisition payback. The gate must also permit the best decision to be no further construction.

\subsection{Why the first mission matters}

The browser begins with an ordinary deliverable: a schedule that respects stated dependencies, pooled workers and calendar limits. Users can inspect the proposed start and finish times, edit assumptions, and export a machine-readable plan. A target miss is displayed as a target miss. A feasible plan is not automatically an optimal one, and a verified model does not certify that its duration estimates describe a real project. This is a usable computational planning instrument, not physical execution of the project.

Later work earns a reason to return only if it improves an outcome the user values. Lower checking burden, a better response to disrupted capacity or an additional feasible deadline can matter; a larger strategy archive need not. Conditional break-even calculations help decide whether to investigate further, but real willingness to pay, accepted delivery and commercial economics require separate observations. No paid customer, token price, revenue, liquidity or profit is inferred from this study.

\subsection{Limits of independence, scope and cost}

The producer, independent checker and adversarial review were developed as separate components within the same commissioned project. The Python checker was authored from the input/output contract without importing the producer's solver. The browser has its own implementation of the acceptance checks. This separation catches specific implementation errors; it is not independent institutional replication, formal verification or a defense against a compromised host administrator. Shared specifications can contain correlated mistakes. The recovered historical artifact audit and the newly completed exact frozen-policy replay are distinct checks, both preserved; neither establishes generalization outside the original benchmark.

The experiment samples a small synthetic scheduling family within an authored strategy language. Its fixtures do not represent the distribution of all release projects. Calendar size, job count, worker demand, grammar and deadlines are bounded; out-of-scope inputs should be refused. A finite stream may not repay an asset that would be useful under a different demand distribution. Conversely, extrapolating break-even without measuring maintenance, stale assumptions and future qualification rates can turn a promising asset into a liability.

The browser-first strategy language contains competent interpreted specialists, while the native mixed-integer optimizer is an additional strong reference. It is not a freely interchangeable action inside the counted-work policy: its opaque native work has a separately bounded timing model. This limits the experiment's claim. A practical hybrid that chooses between fast heuristics and native optimization under a common measured time-and-quality budget could be the better service. A winning specialist should be adopted on that evidence; an agentic label supplies no reason to exclude it. This edition does not claim that its controller is already the best such hybrid.

Declared work units are explicit reference operations, not energy, money, total CPU instructions or equal elapsed time. Native optimization uses a different execution model and is compared with separately disclosed wall and CPU measurements. Setup, research design, authoring, debugging and review are documented in an intervention ledger, but their human durations and model charges are unavailable. Those unknowns prevent an all-in monetary productivity claim. The bounded automatic run has its own exact interventions and resource ledger; it does not erase the work required to build the experiment.

\begin{insight}[The standard for the next successor]
Retain a method because it changes fresh, verified work for the better. Retain the incumbent when it does not. Let the same test decide whether searching for the next method is worth doing.
\end{insight}
